%!TEX root = ../sztuthesis_main.tex
% 设置中文摘要

\phantomsection
\addcontentsline{toc}{section}{摘要}

\keywordscn{预训练与微调；蛋白质-配体结合预测；亲和力预测；姿态排序；分子对接评测}
\categorycn{TP391}
\begin{abstractcn}
蛋白质-配体结合预测是计算机辅助药物设计中的基础问题，其核心目标是在有限计算预算下对蛋白质口袋与候选小分子的结合强度和候选构象质量做出有效评估。现有深度学习方法通常分别围绕亲和力预测或姿态排序展开建模，如何在统一表示学习框架下获得稳定的复合物表征，并进一步适配到具体预测任务，仍然是结构基础药物设计中的重要问题。针对这一需求，本文基于 Graph-Transformer 主干，构建了一套面向蛋白质-配体复合物建模的预训练与微调两阶段方法流程。

本文的主要工作包括：（1）围绕蛋白质-配体复合物表示学习，构建基于 Graph-Transformer 的建模框架，并面向亲和力预测与候选姿态排序任务整理相应的训练与推理流程；（2）采用预训练与微调的两阶段训练策略，先学习稳定的复合物表示，再针对目标任务进行参数适配，以减轻直接从头训练带来的优化不稳定问题；（3）围绕数据预处理、模型训练、checkpoint 选择与离线评测，整理形成较完整的实验流程，并以 Pearson 相关系数、AUROC 和 Top-k success rate 等指标对模型进行评估。

实验围绕亲和力预测与候选姿态排序两个层面展开。结果表明，两阶段训练流程能够支持相关任务的稳定建模与评测，训练和测试阶段的数据处理、模型选择与指标统计过程也得到了系统整理。本文工作为后续围绕蛋白质-配体打分与候选构象重排的模型改进提供了可继续扩展的实验基础。

\end{abstractcn}
